% GENERATED FILE. Edit canonical lessons, then run npm run generate:books.
% canonical-source-hash: fnv1a64:cf34a7dc7c1da301
% canonical-lessons: GU-C114-bajri, GU-C114-ghau, GU-C114-makai, GU-C114-rigan, GU-C114-kobi

\chapter{Millet, Wheat, Maize, Aubergine, Cabbage}
\label{ch:gu-millet-wheat-maize-aubergine-cabbage}
\hlchaptermodality{114}

\begin{chapteropening}
\textbf{By the end of this chapter:} \emph{I can say millet, wheat, maize, an aubergine and cabbage.}

Complete the last lesson of chapter 114: i can say millet, wheat, maize, an aubergine and cabbage.
\end{chapteropening}

\section[bājrī]{\gu{બાજરી} (bājrī) — millet}

\label{lesson:GU-C114-bajri}



\begin{quote}
Before the new one: say the Gujarati for chickpeas, then the Gujarati for mung beans.
\end{quote}

\subsection*{You'll want to know: \gu{બાજરી}}
\textbf{\gu{બાજરી}} (\emph{bājrī}) — \textquotedblleft{}millet\textquotedblright{}.

\begin{grammarlens}[title={What you've built}]
One more word for this chapter.
\end{grammarlens}

\subsection*{Your turn}
\begin{itemize}
\raggedright
  \item \emph{Say it:} \emph{bājrī}
  \item \emph{Say it:} \emph{bājrī}, once more
  \item \emph{Say it:} say \emph{bājrī}
  \item \emph{Recall:} say the Gujarati for chickpeas, then the Gujarati for mung beans, then say \emph{bājrī} again
\end{itemize}

\begin{tcolorbox}[breakable,colback=teal!4,colframe=teal!35!black,title={Before you move on}]
What does \gu{બાજરી} mean? (\textquotedblleft{}Millet\textquotedblright{}.) Say it once more.
\end{tcolorbox}
\section[ghaũ]{\gu{ઘઉં} (ghaũ) — wheat}

\label{lesson:GU-C114-ghau}



\begin{quote}
Before the new one: say the Gujarati for mung beans, then the Gujarati for millet.
\end{quote}

\subsection*{You'll want to know: \gu{ઘઉં}}
\textbf{\gu{ઘઉં}} (\emph{ghaũ}) — \textquotedblleft{}wheat\textquotedblright{}.

\begin{grammarlens}[title={What you've built}]
One more word for this chapter.
\end{grammarlens}

\subsection*{Your turn}
\begin{itemize}
\raggedright
  \item \emph{Say it:} \emph{ghaũ}
  \item \emph{Say it:} \emph{ghaũ}, once more
  \item \emph{Say it:} say \emph{ghaũ}
  \item \emph{Recall:} say the Gujarati for mung beans, then the Gujarati for millet, then say \emph{ghaũ} again
\end{itemize}

\begin{tcolorbox}[breakable,colback=teal!4,colframe=teal!35!black,title={Before you move on}]
What does \gu{ઘઉં} mean? (\textquotedblleft{}Wheat\textquotedblright{}.) Say it once more.
\end{tcolorbox}
\section[makāī]{\gu{મકાઈ} (makāī) — maize}

\label{lesson:GU-C114-makai}



\begin{quote}
Before the new one: say the Gujarati for millet, then the Gujarati for wheat.
\end{quote}

\subsection*{You'll want to know: \gu{મકાઈ}}
\textbf{\gu{મકાઈ}} (\emph{makāī}) — \textquotedblleft{}maize\textquotedblright{}.

\begin{grammarlens}[title={What you've built}]
One more word for this chapter.
\end{grammarlens}

\subsection*{Your turn}
\begin{itemize}
\raggedright
  \item \emph{Say it:} \emph{makāī}
  \item \emph{Say it:} \emph{makāī}, once more
  \item \emph{Say it:} say \emph{makāī}
  \item \emph{Recall:} say the Gujarati for millet, then the Gujarati for wheat, then say \emph{makāī} again
\end{itemize}

\begin{tcolorbox}[breakable,colback=teal!4,colframe=teal!35!black,title={Before you move on}]
What does \gu{મકાઈ} mean? (\textquotedblleft{}Maize\textquotedblright{}.) Say it once more.
\end{tcolorbox}
\section[rī̃gaṇ]{\gu{રીંગણ} (rī̃gaṇ) — an aubergine}

\label{lesson:GU-C114-rigan}



\begin{quote}
Before the new one: say the Gujarati for wheat, then the Gujarati for maize.
\end{quote}

\subsection*{You'll want to know: \gu{રીંગણ}}
\textbf{\gu{રીંગણ}} (\emph{rī̃gaṇ}) — \textquotedblleft{}an aubergine\textquotedblright{}.

\begin{grammarlens}[title={What you've built}]
One more word for this chapter.
\end{grammarlens}

\subsection*{Your turn}
\begin{itemize}
\raggedright
  \item \emph{Say it:} \emph{rī̃gaṇ}
  \item \emph{Say it:} \emph{rī̃gaṇ}, once more
  \item \emph{Say it:} say \emph{rī̃gaṇ}
  \item \emph{Recall:} say the Gujarati for wheat, then the Gujarati for maize, then say \emph{rī̃gaṇ} again
\end{itemize}

\begin{tcolorbox}[breakable,colback=teal!4,colframe=teal!35!black,title={Before you move on}]
What does \gu{રીંગણ} mean? (\textquotedblleft{}An aubergine\textquotedblright{}.) Say it once more.
\end{tcolorbox}
\section[kobī]{\gu{કોબી} (kobī) — cabbage}

\label{lesson:GU-C114-kobi}



\begin{quote}
Before the new one: say the Gujarati for maize, then the Gujarati for an aubergine.
\end{quote}

\subsection*{You'll want to know: \gu{કોબી}}
\textbf{\gu{કોબી}} (\emph{kobī}) — \textquotedblleft{}cabbage\textquotedblright{}.

\begin{grammarlens}[title={What you've built}]
One more word for this chapter.
\end{grammarlens}

\subsection*{Your turn}
\begin{itemize}
\raggedright
  \item \emph{Say it:} \emph{kobī}
  \item \emph{Say it:} \emph{kobī}, once more
  \item \emph{Say it:} say \emph{kobī}
  \item \emph{Recall:} say the Gujarati for maize, then the Gujarati for an aubergine, then say \emph{kobī} again
\end{itemize}

\begin{tcolorbox}[breakable,colback=teal!4,colframe=teal!35!black,title={Before you move on}]
What does \gu{કોબી} mean? (\textquotedblleft{}Cabbage\textquotedblright{}.) Say it once more.
\end{tcolorbox}
